This exam paper investigates how three technologies - Kubernetes, AWS Lambda and Private 5G — relate to edge computing in the specific context of the architectural patterns described in our own Edge Computing report \cite{edge_computing_report}. Edge computing shifts computation toward the physical environment, but this shift exposes structural requirements that traditional cloud-native systems do not address on their own. The three selected topics were not chosen for thematic diversity; each corresponds to a concrete architectural dependency that emerges whenever computation is distributed across numerous, heterogeneous and geographically dispersed nodes.

Kubernetes relates to the operational discipline required in such environments. Our edge prototype already showed that even modest multi-node setups involve configuration management, deployment consistency and workload placement challenges. The Kubernetes report \cite{kubernetes_report} demonstrates how a declarative control plane and continuous reconciliation provide the mechanisms required to maintain coherence across distributed sites. Kubernetes is essential not merely as a typical orchestration solution, but as the indispensable control plane that prevents large-scale edge environments from becoming operationally fragmented.

AWS Lambda addresses a different but equally structural requirement: scalable cloud-side processing for workloads that cannot or should not be executed at the edge. Our own pipeline offloaded anomaly summaries to the cloud, reflecting a broader pattern in edge systems: local processing reduces latency and bandwidth, but global aggregation, enrichment and asynchronous workflows still require elastic execution environments. Lambda’s event-driven model offers a principled mechanism to integrate cloud logic into edge architectures without maintaining persistent infrastructure \cite{lambda_report}.

Private 5G is relevant because edge computing depends on predictable connectivity, and conventional wireless networks cannot provide the latency bounds, mobility guarantees or isolation levels that many real-world edge workloads require. The Private 5G report \cite{private_5g_report} shows how licensed-spectrum control, deterministic scheduling and local user-plane placement fundamentally change the performance envelope in which edge systems operate. This connection becomes most explicit in the final chapter through Aether, which integrates SD-Core, SD-RAN, MEC workloads and Kubernetes orchestration into a unified cloud-native system. Aether demonstrates that network behaviour and compute behaviour cannot be considered independently in distributed architectures: the viability of edge computing depends on deterministic communication, while the usefulness of Private 5G becomes most apparent when paired with a structured edge execution layer. 

The relationships studied in the following chapters therefore reflect architectural necessity rather than thematic convenience. Each technology addresses a specific limitation that arises when computation moves toward the edge, and their interplay becomes clearest in systems, like Aether, where orchestration, connectivity and distributed compute converge into a single operational model.
